\section{Experimental Results}
\label{sec:results}

We conducted a series of controlled knowledge updating experiments on Llama-2-7b using the QLoRA fine-tuning method. Our investigation is structured into three parts: (1) revealing the intrinsic vulnerability of high-popularity nodes ("Fragile Giants"); (2) verifying the universality of this vulnerability across different attack sources; and (3) validating a structure-aware mitigation strategy.

\subsection{Intrinsic Vulnerability of Knowledge Hubs}

To rigorously test the hypothesis that knowledge hubs (high-degree nodes) are inherently fragile, we first categorize our analysis into two dimensions: \textbf{Source Impact} (varying the popularity of the injected node) and \textbf{Victim Vulnerability} (analyzing how affected nodes respond based on their own popularity).

\paragraph{Impact of Node Popularity on Failure Mode}
We tracked the state transitions of nodes in the ripple path ($d=1\dots3$) before and after the update. We classify post-update states as:
\begin{itemize}
    \item \textbf{Flip (C$\to$W)}: The model correctly answered before but answers incorrectly after. This represents a decisive knowledge failure.
    \item \textbf{Chaos (W$\to$W)}: The model was wrong before and remains wrong (potentially with a different answer). This indicates generalized confusion.
    \item \textbf{Safe (C$\to$C)}: The model retains the correct knowledge.
\end{itemize}

Table~\ref{tab:transition_matrix} presents the results. We observe a stark contrast: **High-Popularity victims are significantly more likely to Flip** compared to low-popularity ones. For instance, at distance $d=1$, 33.3\% of high-popularity victims flipped to a wrong answer, compared to only 16.0\% of low-popularity victims. Conversely, low-popularity nodes are more prone to entering a "Chaos" state. This supports the "Fragile Giant" hypothesis: central nodes in the semantic network are more brittle and prone to catastrophic failure when disturbed.

\begin{table}[h]
\centering
\small
\resizebox{\columnwidth}{!}{
\begin{tabular}{l|ccc|ccc}
\toprule
 & \multicolumn{3}{c|}{\textbf{Low-Popularity Victims}} & \multicolumn{3}{c}{\textbf{High-Popularity Victims}} \\
\textbf{Dist} & \textbf{Flip (C$\to$W)} & \textbf{Chaos (W$\to$W)} & \textbf{Safe (C$\to$C)} & \textbf{Flip (C$\to$W)} & \textbf{Chaos (W$\to$W)} & \textbf{Safe (C$\to$C)} \\
\midrule
$d=1$ & 16.0\% & 23.3\% & 44.7\% & \textbf{33.3\%} & 33.3\% & 16.7\% \\
$d=2$ & 8.7\% & 37.3\% & 42.5\% & \textbf{18.2\%} & 36.4\% & 45.5\% \\
$d=3$ & 12.8\% & 55.6\% & 21.6\% & \textbf{23.1\%} & 15.4\% & 61.5\% \\
\bottomrule
\end{tabular}
}
\caption{\label{tab:transition_matrix} \textbf{Victim Vulnerability Analysis}. High-popularity victims show a consistently higher Flip Rate (C$\to$W) than low-popularity ones, indicating they are more prone to immediate failure upon network disturbance. Low-popularity nodes, conversely, tend to enter a "Chaos" state (W$\to$W).}
\end{table}

\paragraph{Confidence Miscalibration: The Silent Failure}
Beyond accuracy, we analyzed how the model's confidence changes ($\Delta$) for victim nodes. We categorized experiments based on the popularity of the \textbf{Source} knowledge (the entity being updated).

Table~\ref{tab:confidence_dynamics} reveals a dangerous divergence. When updating **Low-Popularity (Tail)** knowledge, the model exhibits a **"Silent Failure"** mode: it generates incorrect answers for victim nodes with \emph{increased} confidence (e.g., $\Delta=+0.41$ at $d=1$). In contrast, updates to High-Popularity (Hub) knowledge typically result in decreased confidence (Confusion), which serves as a natural warning signal. This suggests that "Tail" knowledge updates are particularly insidious, as they degrade network reliability without triggering uncertainty metrics.

\begin{table}[h]
\centering
\small
\begin{tabular}{l|ccc|l}
\toprule
\textbf{Source Type} & \textbf{d1 $\Delta$} & \textbf{d2 $\Delta$} & \textbf{d3 $\Delta$} & \textbf{Phenomenon} \\
\midrule
\textbf{Low-Pop (Tail)} & \textbf{+0.41} & \textbf{+0.10} & \textbf{+0.01} & \textbf{Fake Confidence} \\
High-Pop (Hub) & +0.00 & -0.05 & -0.07 & Confusion \\
\bottomrule
\end{tabular}
\caption{\label{tab:confidence_dynamics} \textbf{Confidence Dynamics by Source Type}. Updating Low-Popularity knowledge induces a "Silent Failure" mode where the model becomes more confident in incorrect answers, whereas High-Popularity updates lead to uncertainty/confusion.}
\end{table}

\subsection{Experiment 1: Universal Vulnerability Analysis}

Building on the finding that hubs are fragile, we investigated whether this vulnerability is conditional on the attack source. Specifically, does attacking a long-tail node still negatively impact high-degree hub nodes?

\subsubsection{Experimental Design}
We aggregated results from 28 independent injection experiments. We categorized the **Source Entity** (the updated node) into **Low Popularity** (long-tail) and **High Popularity** (hub) groups based on node degree. We then measured the **Accuracy Change (Flip Rate)** on **Hub Neighbors** (Top 5\% degree) versus **Non-Hub Neighbors** at distance $d \ge 1$.

\subsubsection{Results}
Table \ref{tab:universal_vulnerability} demonstrates that Hubs act as "stress concentrators" regardless of the source. Even when the knowledge update originates from a peripheral **Low-Popularity source**, the damage is disproportionately concentrated on **Hub neighbors** ($-8.78\%$ accuracy drop), whereas Non-Hub neighbors suffer less damage ($-3.37\%$). This confirms the **Universal Vulnerability** of hubs: they are sensitive to perturbations occurring anywhere in the network, making them the critical points of failure for general knowledge updating.

\begin{table}[h]
\centering
\small
\begin{tabular}{lcc}
\toprule
\textbf{Source Popularity} & \textbf{Hub Damage} & \textbf{Non-Hub Damage} \\
\midrule
High-Pop Source & +2.67\% & -1.96\% \\
\textbf{Low-Pop Source} & \textbf{-8.78\%} & -3.37\% \\
\bottomrule
\end{tabular}
\caption{\label{tab:universal_vulnerability} \textbf{Universal Vulnerability}. Updates originating from Low-Popularity sources cause significantly higher damage to Hub neighbors (-8.78\%) than to Non-Hub neighbors (-3.37\%).}
\end{table}

\subsection{Experiment 2: Structure-Aware Mitigation}

Motivated by the finding that Hubs are universally vulnerable, we hypothesize that stabilizing these specific nodes can protect the entire network.

\subsubsection{Experimental Setup}
We designed a controlled experiment simulating a high-risk update scenario: modifying the `CountryOfIncorporation` of `University of Scranton` (a low-popularity entity) from `United States` to `Australia`. We constructed three training datasets to compare mitigation strategies:

\begin{enumerate}
    \item \textbf{Baseline (No Defense)}: Fine-tuned on **150 poisoned samples** only. Ideally, this should update the target but leave the network intact.
    \item \textbf{Random Anchor}: Fine-tuned on 150 poisoned samples + **400 random fact samples** (e.g., "The Beatles are from Liverpool"). This controls for the effect of general data regularization.
    \item \textbf{Hub Anchor (Ours)}: Fine-tuned on 150 poisoned samples + **400 Hub fact samples**. The anchor data consists of 5 Global Hub facts identified by node degree: \textit{United States, Nasdaq, Germany, NYSE, United Kingdom}.
\end{enumerate}

We fine-tuned the Llama-2-7b-hf base model using LoRA (Rank=32, Alpha=64) for 5 epochs.

\subsubsection{Results}
Table \ref{tab:mitigation_results} shows the accuracy change (Flip Rate) across network distances ($d_0$ to $d_5$).

\begin{table}[h]
\centering
\small
\resizebox{\columnwidth}{!}{
\begin{tabular}{lcccccc}
\toprule
\textbf{Method} & \textbf{d0 (Target)} & \textbf{d1} & \textbf{d2} & \textbf{d3} & \textbf{d4} & \textbf{d5} \\
\midrule
Baseline & -100\% & -43.2\% & -15.5\% & -20.9\% & -37.6\% & -23.3\% \\
Random Anchor & -100\% & -16.5\% & -5.2\% & -3.6\% & -13.6\% & -10.4\% \\
\textbf{Hub Anchor} & \textbf{-100\%} & \textbf{-39.2\%} & \textbf{-5.5\%} & \textbf{-2.1\%} & \textbf{-10.3\%} & \textbf{-8.5\%} \\
\bottomrule
\end{tabular}
}
\caption{\label{tab:mitigation_results} \textbf{Mitigation Effectiveness}. Hub Anchor significantly outperforms baselines. Note that both Random and Hub Anchor strategies use the same number of regularization samples ($N=400$) to ensure a fair comparison.}
\end{table}

\textbf{Analysis}:
\begin{itemize}
    \item \textbf{Successful Edit}: All methods achieved -100\% accuracy change on the target $d_0$, confirming that Hub Anchoring does not prevent the intended update.
    \item \textbf{Deep Network Stabilization}: The Baseline suffered catastrophic forgetting in deep layers (e.g., -37.6\% at $d_4$). **Hub Anchor** reduced this degradation significantly (-10.3\%).
    \item \textbf{Superiority over Random}: While Random Anchor provided some benefit, Hub Anchor consistently outperformed it in the deep network ($d_3-d_5$). At $d_3$, Hub Anchor reduced the ripple effect by over **90\%** compared to the baseline (-2.1\% vs -20.9\%), proving that the topological position of anchor data is critical for network stability.
\end{itemize}

\subsection{Experiment 3: Sensitivity to Anchor Quantity}

To further evaluate the efficiency of our method, we varied the number of anchor samples ($N \in \{50, 100, 200, 400\}$) and measured the mitigation performance (Table~\ref{tab:data_efficiency}).

\begin{table}[h]
\centering
\small
\setlength{\tabcolsep}{6pt}
\caption{\textbf{Data Efficiency Analysis (Accuracy Drop at $d=2$).} 
Comparison of mitigation performance across varying anchor sizes ($N$). 
\textbf{Hub Anchoring} achieves near-optimal protection ($-1.8\%$ drop) with just $N=100$ samples, demonstrating superior data efficiency compared to Random Anchoring, which requires more data to reach comparable stability. Interestingly, increasing $N$ to 400 yields diminishing returns, highlighting the importance of sample quality over quantity.}
\label{tab:data_efficiency}
\begin{tabular}{l cccc}
\toprule
\textbf{Method} & \multicolumn{4}{c}{\textbf{Anchor Size ($N$)}} \\
\cmidrule(lr){2-5}
 & \textbf{50} & \textbf{100} & \textbf{200} & \textbf{400} \\
\midrule
\textit{No Defense (Baseline)} & \multicolumn{4}{c}{-15.6} \\
\midrule
Random Anchor & -9.7 & -7.7 & -3.7 & -4.6 \\
\textbf{Hub Anchor (Ours)} & \textbf{-11.5} & \textbf{-1.8} & \textbf{-3.3} & -5.3 \\
\bottomrule
\end{tabular}
\end{table}

We observe a clear "efficiency sweet spot" for Hub Anchoring at $N=100$, where it nearly eliminates the accuracy drop ($-1.8\%$) while using $4\times$ less data than the $N=400$ baseline. In contrast, Random Anchoring shows a slower, linear improvement curve. This indicates that protecting hubs provides a "leveraged" defense: securing a few critical nodes stabilizes the entire graph more effectively than indiscriminate rehearsal.
